%%%%%%%%%%%%%%%%%%%%%%%%%%%%%%%%%%%%%%%%%%%%%%%%%
%%%%%%%%%%%% chap: implementation %%%%%%%%%%%%%%%%%
%%%%%%%%%%%%%%%%%%%%%%%%%%%%%%%%%%%%%%%%%%%%%%%%%

\chapter{System Design and Implementation}\label{chapter:chap3}

\section{System Architecture Overview}

The proposed system is a multi-stage phishing honeypot platform designed to
capture and analyse the complete behavioural lifecycle of a phishing attack,
from initial credential harvesting through post-compromise SSH activity. The
architecture consists of six integrated components deployed on a single Ubuntu
host: a Flask-based phishing decoy, a Cowrie SSH honeypot, a PostgreSQL
credential store, an automated bot simulation framework, an ELK-based log
pipeline, and a machine learning classification module.

The central design principle is the use of a shared \texttt{session\_id}
identifier that links events across all stages of an attack. When a simulated
victim submits credentials to the phishing decoy, a session identifier is
generated and persisted to PostgreSQL. The attacker bot subsequently reads this
identifier and uses it when initiating an SSH session against Cowrie, ensuring
that phishing events and post-compromise SSH events can be correlated as a
single end-to-end behavioural record. This cross-stage linkage is a key
contribution of the proposed system, as existing phishing honeypots typically
capture only the credential submission stage without tracking subsequent
attacker activity \cite{Leita2008Phishing,Bringer2012HoneypotSurvey}.

The system is deployed using Docker Compose, with the exception of Cowrie,
which runs directly on the host operating system to avoid SSH port conflicts
inherent to containerised deployments. Each of the four bot profiles is
assigned to a dedicated Docker network subnet, resulting in genuinely distinct
source IP addresses per attacker profile. This design ensures that the machine
learning pipeline must learn from behavioural features rather than trivially
classifying sessions by IP address alone.

\section{Phishing Decoy Infrastructure}

The phishing decoy is implemented as a three-stage Flask web application
impersonating a generic corporate login portal named CorpNet. The three stages
correspond to the principal phases of a modern phishing campaign: credential
harvesting, multi-factor authentication bypass, and post-authentication
redirection \cite{Leita2008Phishing}.

\textbf{Stage 1 — Credential harvesting.} The login page presents a realistic
corporate interface soliciting an email address and password. In addition to
the submitted credentials, the application captures a set of behavioural
indicators using client-side JavaScript: the time elapsed between page load and
form submission (\texttt{time\_on\_page\_ms}), the delay between the first
keystroke and form submission (\texttt{typing\_delay\_ms}), and the number of
field corrections made during input (\texttt{field\_corrections}). These three
features capture timing and interaction patterns that differ significantly
between human users and automated bots, and are used as primary features in the
machine learning pipeline described in Section~\ref{sec:ml_design}.

\textbf{Stage 2 — MFA bypass.} Upon credential submission, the application
generates a six-digit one-time code and delivers it via SMTP to a locally
deployed MailHog instance, which functions as a contained fake mail server. The
victim bot retrieves this code programmatically from the MailHog REST API and
submits it to the MFA endpoint. Real phishing infrastructure routinely
implements MFA relay to bypass second-factor controls; this stage replicates
that behaviour in a controlled and ethical manner \cite{Moric2025Honeypots}.

\textbf{Stage 3 — Redirection.} Successful MFA submission redirects the
simulated victim to a fake dashboard page, completing the session. The
application logs a \texttt{phishing.session.complete} event containing the full
session duration and username, which closes the phishing record in PostgreSQL.

All events are written as structured JSON to a log file mounted into the
Filebeat container, which forwards them to Elasticsearch under the index pattern
\texttt{honeypot-flask-*}. Events are simultaneously persisted to PostgreSQL
via the \texttt{phishing\_sessions} table for use by the machine learning
pipeline.

\section{Cowrie SSH Honeypot}
\label{sec:cowrie}

The third stage of the attacker lifecycle is post-credential exploitation via
SSH\@. To capture this stage, the system deploys Cowrie \cite{Song2005HoneyStat},
a medium-interaction SSH honeypot that emulates a full Unix shell environment
without exposing a real operating system.

\subsection{Deployment Architecture}

Cowrie runs directly on the host machine rather than inside a Docker container.
This design decision avoids a port-forwarding conflict: Docker's SSH daemon
occupies port~22, so placing Cowrie in a container would require a second level
of port mapping that complicates the bot-to-honeypot connection. Running bare on
the host allows Cowrie to bind directly to port~2222 and be reached by all four
bot containers via the host gateway address.

The honeypot is configured in \emph{shell} backend mode, meaning all commands
are handled by Cowrie's built-in command emulator rather than forwarded to a
real system. This is appropriate for the thesis scenario because the goal is
behavioural logging, not high-fidelity command execution. The relevant
configuration parameters are shown in Table~\ref{tab:cowrie_cfg}.

\begin{table}[h]
\centering
\caption{Key Cowrie configuration parameters.}
\label{tab:cowrie_cfg}
\begin{tabular}{ll}
\hline
\textbf{Parameter} & \textbf{Value} \\
\hline
Hostname              & \texttt{CORP-SRV-01} \\
Backend               & \texttt{shell} \\
Authentication        & \texttt{UserDB} \\
Listen port           & 2222 \\
Idle timeout          & 180 s \\
Authentication timeout & 120 s \\
TTY logging           & Enabled \\
Log format            & Rotating JSON \\
\hline
\end{tabular}
\end{table}

The hostname \texttt{CORP-SRV-01} reinforces the corporate deception established
by the phishing page. An attacker who steals credentials from what appears to be
a corporate login portal and subsequently reaches a server by this name is more
likely to proceed with reconnaissance, increasing the volume and quality of
captured behaviour \cite{Holz2006AttackPatterns}.

\subsection{Credential Configuration}

Authentication is controlled by \texttt{userdb.txt}, which lists every
username–password pair that Cowrie will accept. The file contains nine entries:
the five planted corporate identities that are also seeded in the Flask decoy
and PostgreSQL (\texttt{john.mitchell}, \texttt{sarah.chen}, \texttt{admin},
\texttt{it.support}, \texttt{j.smith}), two common default credentials
(\texttt{root/123456}, \texttt{admin/admin}), and a wildcard entry that accepts
any remaining credential pair. The wildcard ensures that scanners probing with
arbitrary credentials are also admitted, generating SSH session data rather than
being silently rejected at the authentication stage.

\subsection{Fake Filesystem}

Cowrie's \texttt{honeyfs/} directory overlays a set of files onto the emulated
filesystem designed to elicit further attacker interaction. Three files are of
particular interest:

\begin{itemize}
    \item \textbf{\texttt{/etc/passwd}} — Contains eleven entries including the
    five planted corporate users with realistic home directories and login
    shells. This file is the first target of account-discovery commands such as
    \texttt{cat /etc/passwd}.

    \item \textbf{\texttt{/etc/secrets.txt}} — A simulated internal
    configuration file containing fictitious database credentials, AWS access
    keys, and SMTP passwords. Planting credentials of this kind follows the
    deception principle of breadcrumbing \cite{Leita2008Phishing}: each secret
    gives the attacker a reason to continue and attempt lateral movement,
    extending the session duration and command count.

    \item \textbf{\texttt{/home/john.mitchell/.bash\_history}} — A
    pre-populated command history including SSH connections to internal
    addresses and database access commands. These entries simulate artefacts
    left by a legitimate administrator and encourage sophisticated attackers to
    investigate internal network topology.
\end{itemize}

\subsection{Logging}

Cowrie produces structured JSON logs at
\texttt{var/log/cowrie/cowrie.json}. Each event records the event type
(\texttt{cowrie.login.success}, \texttt{cowrie.command.input},
\texttt{cowrie.session.closed}), the source IP, session identifier, timestamp,
and — for command events — the exact string entered by the attacker. TTY
logging is also enabled, producing a full terminal replay of every session.
The \texttt{session\_id} written into SSH sessions by the attacker bot is the
same UUID generated during the Flask phishing phase, enabling cross-stage
correlation described in Section~\ref{sec:ml_design}.

\section{Attacker Bot Simulation}
\label{sec:bots}

Evaluating a honeypot in a controlled laboratory environment requires a source
of realistic attacker traffic. Four automated bot profiles are implemented in
\texttt{bot\_manager.py}, each parameterised from real-world data and running
in its own isolated Docker container and network subnet.

\subsection{Justification for Synthetic Traffic}

The use of synthetic traffic in honeypot research is established practice when
live exposure is not feasible \cite{Kreibich2005Replay,Alata2006Honeynet}. The
behavioural parameters of each profile are derived from two authoritative
sources:

\begin{enumerate}
    \item \textbf{CyberLab Honeynet Dataset} \cite{Alata2006Honeynet} — 293
    real SSH honeypot sessions from which login timing distributions, session
    duration medians, and login attempt counts are extracted.
    \item \textbf{MITRE ATT\&CK} — Techniques T1078 (Valid Accounts), T1059
    (Command and Scripting Interpreter), T1087 (Account Discovery), and T1005
    (Data from Local System) are used to select the command sequences executed
    by each profile.
\end{enumerate}

This dual grounding means that while the traffic is synthetic, the distributions
it samples from are empirically validated.

\subsection{Two-Phase Attack Model}

Each bot cycle executes two phases that mirror the complete attacker lifecycle:

\begin{description}
    \item[Phase 1 — Victim bot] Simulates a user who receives a phishing email,
    visits the Flask decoy, submits credentials, and completes MFA\@. The bot
    reads the MFA code from the MailHog API, replicating the account takeover
    step without any external network dependency.

    \item[Phase 2 — Attacker bot] After a profile-specific delay, reads the
    stolen credentials from PostgreSQL and opens an SSH connection to Cowrie on
    port~2222. It then executes the command sequence defined for its profile.
\end{description}

Both phases share the same \texttt{session\_id} generated at the start of
Phase~1. This identifier is written into \texttt{phishing\_sessions},
\texttt{ssh\_sessions}, and \texttt{ml\_features}, enabling every event in the
attacker lifecycle to be traced to a single origin.

\subsection{Profile Definitions}

The four profiles and their key parameters are summarised in
Table~\ref{tab:bot_profiles}.

\begin{table}[h]
\centering
\caption{Attacker bot profile parameters derived from CyberLab distributions and MITRE ATT\&CK.}
\label{tab:bot_profiles}
\begin{tabular}{lllll}
\hline
\textbf{Profile} & \textbf{Login delay} & \textbf{MFA delay} & \textbf{Commands} & \textbf{MITRE} \\
\hline
Scanner            & 0.1–0.5 s  & 0.1–0.3 s   & 4  & T1087 \\
Credential stuffer & 0.5–2.0 s  & 0.2–0.5 s   & 4  & T1110 \\
Human-like         & 3.0–12.0 s & 8.0–30.0 s  & 11 & T1059, T1005 \\
Sophisticated      & 8.0–25.0 s & 15.0–45.0 s & 16 & T1078, T1087, T1005 \\
\hline
\end{tabular}
\end{table}

\paragraph{Scanner.}
Parameterised from 34 CyberLab sessions with a median duration of 25.4~s.
The scanner submits all five credential pairs in rapid succession, uses
non-browser user agents (\texttt{python-requests}, \texttt{curl},
\texttt{Nikto}), and executes four basic account-discovery commands with
inter-command delays of 0.1–0.5~s. This profile represents automated
vulnerability scanning rather than targeted intrusion
\cite{Bringer2012HoneypotSurvey}.

\paragraph{Credential stuffer.}
Parameterised from 4 CyberLab sessions with a mean of 23.75 login attempts.
It submits ten credential pairs at 0.5–2.0~s per field, consistent with
list-based brute-force tooling. The command set targets privilege escalation
indicators (\texttt{cat /etc/shadow}, \texttt{ls /home})
\cite{Bailey2009Botnets}.

\paragraph{Human-like.}
Parameterised from 148 CyberLab sessions with a median session duration of
52.1~s. Login delays of 3–12~s and MFA delays of 8–30~s reflect the timing of
a human operator reading and typing. The command sequence of eleven entries
includes network enumeration and data collection, consistent with MITRE ATT\&CK
T1059 and T1005 \cite{Holz2006AttackPatterns}.

\paragraph{Sophisticated.}
Represents an Advanced Persistent Threat actor. Login delays of 8–25~s and
MFA delays of 15–45~s indicate deliberate, unhurried operation. The
sixteen-command sequence includes persistent access investigation
(\texttt{crontab -l}), broad configuration file discovery, and lateral movement
indicators. The profile executes a single login attempt and uses only a
legitimate browser user agent, minimising the automated footprint visible to
anomaly detectors \cite{Freiling2010Evasion}.

\subsection{Network Isolation}

Each bot runs in a dedicated Docker subnet (Table~\ref{tab:bot_subnets}),
ensuring that the source IP address recorded in every session genuinely differs
per profile. This prevents the classifier from memorising IP prefixes rather
than learning from behavioural features \cite{Rajab2007Botnet}.

\begin{table}[h]
\centering
\caption{Docker subnet assignments for bot isolation.}
\label{tab:bot_subnets}
\begin{tabular}{lll}
\hline
\textbf{Profile} & \textbf{Subnet} & \textbf{Bot IP} \\
\hline
Scanner            & \texttt{172.20.0.0/24} & \texttt{172.20.0.5} \\
Credential stuffer & \texttt{172.21.0.0/24} & \texttt{172.21.0.5} \\
Human-like         & \texttt{172.22.0.0/24} & \texttt{172.22.0.5} \\
Sophisticated      & \texttt{172.23.0.0/24} & \texttt{172.23.0.5} \\
\hline
\end{tabular}
\end{table}

\section{Log Collection and Visualisation}
\label{sec:elk}

Centralised log aggregation is implemented using the Elastic Stack
\cite{Holz2006AttackPatterns}: Filebeat ships logs from both honeypot
components to Elasticsearch, where they are indexed and made available for
analysis in Kibana.

\subsection{Filebeat Configuration}

Filebeat is deployed as a Docker service with bind-mount access to both log
directories. It is configured with two inputs and a single Elasticsearch
output, as shown in Table~\ref{tab:filebeat}.

\begin{table}[h]
\centering
\caption{Filebeat input and index configuration.}
\label{tab:filebeat}
\begin{tabular}{lll}
\hline
\textbf{Source} & \textbf{Log path} & \textbf{Elasticsearch index} \\
\hline
Flask phishing decoy & \texttt{/logs/flask/*.json}       & \texttt{honeypot-flask-\{date\}} \\
Cowrie SSH honeypot  & \texttt{/logs/cowrie/cowrie.json} & \texttt{honeypot-cowrie-\{date\}} \\
\hline
\end{tabular}
\end{table}

Both inputs use \texttt{json.keys\_under\_root: true}, which promotes all JSON
fields to the top level of the Elasticsearch document, simplifying Kibana query
construction. Each document is tagged with a \texttt{log\_source} field and a
\texttt{honeypot\_component} field, enabling per-component filtering within a
single Kibana data view. Index lifecycle management is disabled and a custom
index template matching \texttt{honeypot-*} is applied to ensure consistent
field mappings across both index families.

\subsection{Elasticsearch and Kibana}

Elasticsearch~8.13.0 runs as a single-node cluster on port~9200 within the
\texttt{honeypot-net} Docker network. The single-node configuration is
sufficient for the thesis scale and eliminates coordination overhead.

Kibana~8.13.0 is accessible on port~5601 and provides real-time dashboards over
the collected honeypot data. A single data view covering the \texttt{honeypot-*}
index pattern is configured, allowing analysts to query both Flask and Cowrie
events in the same interface and correlate them by \texttt{session\_id} or
timestamp.

\section{Machine Learning Pipeline}
\label{sec:ml_design}

The ML pipeline classifies each recorded session along two independent axes,
constituting a multi-label classification problem \cite{Moric2025Honeypots}.

\subsection{Problem Formulation}

Each session in \texttt{ml\_features} is assigned two ground-truth labels set
by the bot manager at collection time:

\begin{itemize}
    \item \textbf{Bot type} $\in$ \{scanner, credential\_stuffer, human\_like,
    sophisticated\} — identifies \emph{who} the attacker is.
    \item \textbf{Attack stage} $\in$ \{cred\_submission, ssh\_attack\} —
    identifies \emph{what phase} of the lifecycle the session belongs to.
\end{itemize}

The two labels are treated independently: one Random Forest classifier is
trained per label, both using the same feature vector.

\subsection{Feature Engineering}

Six features are extracted for each session, spanning both the phishing and SSH
stages of the lifecycle (Table~\ref{tab:features}). For credential submission
rows the SSH features are set to zero; for SSH attack rows the phishing features
are set to zero. This encoding is deliberate: the zero values are themselves a
discriminating signal that the classifier can exploit.

\begin{table}[h]
\centering
\caption{ML feature set used for both classifiers.}
\label{tab:features}
\begin{tabular}{lll}
\hline
\textbf{Feature} & \textbf{Stage} & \textbf{Description} \\
\hline
\texttt{time\_on\_page\_ms}        & Phishing & Time from page load to form submit \\
\texttt{typing\_delay\_ms}         & Phishing & Mean delay between keystrokes \\
\texttt{field\_corrections}        & Phishing & Number of backspace/correction events \\
\texttt{command\_count}            & SSH      & Total commands executed in session \\
\texttt{session\_duration\_s}      & SSH      & Total SSH session duration (seconds) \\
\texttt{inter\_cmd\_delay\_avg\_s} & SSH      & Mean delay between consecutive commands \\
\hline
\end{tabular}
\end{table}

\subsection{Classifier Design}

Both classifiers use Random Forest \cite{Bringer2012HoneypotSurvey} with 200
estimators, no depth limit, and \texttt{class\_weight=balanced}. Sample weights
are inversely proportional to class frequency, compensating for the imbalance
between the large scanner and credential stuffer classes and the smaller
human-like and sophisticated classes without discarding data.

\subsection{Evaluation Protocol}

Model performance is estimated using 5-fold stratified cross-validation, which
preserves the class distribution in each fold. Three metrics are reported:
accuracy, weighted F1, and macro F1. The macro F1 is the primary evaluation
metric because it treats all classes equally, penalising poor performance on
minority classes — the most forensically interesting attacker profiles.

\subsection{Baseline Comparison}

To contextualise the results, performance is compared against published
classification results on the CICIDS2017 dataset \cite{Bringer2012HoneypotSurvey},
the standard benchmark for network intrusion detection. CICIDS2017 contains
approximately 500\,000 labelled network flows across 14 attack categories.
Reported macro F1 scores on CICIDS2017 for Random Forest range from 0.89 to
0.97 depending on the feature set applied. The honeypot dataset presents a
harder problem in one respect — four-class behavioural profiling rather than
binary anomaly detection — and a more tractable one in another: ground truth is
exact because labels are set by the generating bots rather than inferred by
analysts.
